\documentclass[12pt]{article}

\usepackage[a4paper, total={6in, 8in}]{geometry}
\usepackage{textcomp}

\begin{document}
\setlength{\parindent}{0pt}

\def \t {\textrightarrow}
\def \v {\vspace{1em}}
\def \bi {\begin{itemize}}
\def \ei {\end{itemize}}
\def \s[#1] {\section*{#1}}
\def \ss[#1] {\subsection*{#1}}
\def \sss[#1] {\subsubsection*{#1}}

\s[Jane Eyre]
\ss[Out there in the world]
Last part is protofeminist \t women should be able to do what they want to do \\
Sometimes she climbs to the rooftop and imagines the busy city \\
A good victorian woman was not supposed to be restless \\
My mind's eye \t faculty of imagination \t romantic background \t also inword ear \\
"Women are supposed to be very calm" \t beginning of protofeminist \t but they are not, because they are just like men

\ss[Bertha Mason - page 56]
Bertha belongs to two minorities: she is a woman and comes from the colony \\
She il locked in the attic \t she is referred to in the novel as "The mad woman in the attic" \\
Bertha Mason is described as a villain \t and is presented like an animal \\
Line 35 \t use of irony for the description of married life \t then compaers Bertha to Jane \\
Night of Egypt \t biblical reference

\v

Jane is:
\bi
    \item outspoken 
    \item endwed with a sense of justice \t she gets furious 
    \item restless \t feels the need to travel
    \item imaginative \t when she climbs to the rooftop, she imagines what the world is really like
    \item active \t not a passive victorian woman
    \item rebellious \t she praises to people that are rebelling
    \item she is also introspective
\ei
\end{document}
